\documentclass[11pt,a4paper]{article}
\usepackage[utf8]{inputenc}
\usepackage[vietnamese]{babel}
\usepackage{amsmath,amssymb}
\usepackage{geometry}
\geometry{a4paper, total={180mm,260mm}, left=15mm, right=15mm, top=18mm, bottom=18mm}
\usepackage{tikz}
\usetikzlibrary{shapes.geometric, arrows.meta, positioning, calc}
\usepackage{tcolorbox}
\tcbuselibrary{skins,breakable}
\usepackage{xcolor}
\usepackage{multicol}
\usepackage{array}
\usepackage{fontawesome}
\usepackage{fancyhdr}

\pagestyle{fancy}
\fancyhf{}
\renewcommand{\headrulewidth}{0.5pt}
\renewcommand{\footrulewidth}{0.5pt}
\fancyhead[L]{\small\bfseries\color{teal}\faBook\ SỔ TAY KIẾN THỨC TOÁN 3}
\fancyhead[R]{\small\bfseries\color{teal}Lớp Toán Thầy Lĩnh -- loptoanthaylinh.vn}
\fancyfoot[C]{\small\color{gray}-- Sổ tay kiến thức: Bài 67. Thực hành xem đồng hồ, xem lịch --}

\definecolor{primary}{HTML}{005F73}
\definecolor{secondary}{HTML}{0A9396}
\definecolor{accent}{HTML}{EE9B00}
\definecolor{danger}{HTML}{AE2012}
\definecolor{boxbg}{HTML}{F4F6F8}

\begin{document}

% --- BANNER TIÊU ĐỀ ---
\begin{tcolorbox}[
  enhanced,
  colback=white,
  colframe=primary,
  arc=3mm,
  boxrule=1.2pt,
  left=4mm, right=4mm, top=2.5mm, bottom=2.5mm,
  drop shadow=black!10!white
]
  \centering
  {\Large\bfseries\color{primary}\faBookmark\ BÀI 67: THỰC HÀNH XEM ĐỒNG HỒ, XEM LỊCH}\\[1.5mm]
  {\small\itshape\color{black!75}Kĩ năng đọc tờ lịch tháng, tính khoảng cách thời gian giữa các ngày trong tuần và quản lý thời gian trên 1 trang A4}
\end{tcolorbox}

\vspace{-2mm}


\noindent
\begin{minipage}[t]{0.49\textwidth}\vspace{0pt}
\begin{tcolorbox}[
  enhanced,
  colback=blue!4!white,
  colframe=blue!75!black,
  arc=2mm, boxrule=0.8pt,
  left=2mm, right=2mm, top=1.8mm, bottom=1.8mm,
  title={\footnotesize\bfseries\color{white}\faKey\ 1. Kĩ năng đọc tờ lịch tháng}
]
{\footnotesize
\begin{itemize}\setlength{\itemsep}{1.5pt}\setlength{\leftskip}{-2.5mm}
  \item Mỗi tờ lịch tháng hiển thị tên tháng, năm và bảng các ngày trong tháng.
  \item \textbf{Các hàng ngang:} Chỉ các tuần lễ.
  \item \textbf{Các cột dọc:} Chỉ các ngày trong tuần (Thứ Hai, Thứ Ba, Thứ Tư, Thứ Năm, Thứ Sáu, Thứ Bảy, Chủ Nhật).
  \item \textbf{Quy tắc chu kì tuần lễ (7 ngày):}\\
  Hai ngày cùng thứ ở hai tuần liên tiếp cách nhau đúng $\mathbf{7\text{ ngày}}$:
  $$\text{Ngày tuần sau} = \text{Ngày tuần này} + 7$$
\end{itemize}
}
\end{tcolorbox}
\end{minipage}
\hfill
\begin{minipage}[t]{0.49\textwidth}\vspace{0pt}
\begin{tcolorbox}[
  enhanced,
  colback=teal!4!white,
  colframe=teal!75!black,
  arc=2mm, boxrule=0.8pt,
  left=2mm, right=2mm, top=1.8mm, bottom=1.8mm,
  title={\footnotesize\bfseries\color{white}\faClockO\ 2. Kĩ năng tính khoảng thời gian}
]
{\footnotesize
\begin{itemize}\setlength{\itemsep}{1.5pt}\setlength{\leftskip}{-2.5mm}
  \item \textbf{Khoảng thời gian trôi qua:}
  $$\text{Thời gian trôi qua} = \text{Giờ kết thúc} - \text{Giờ bắt đầu}$$
  \item \textbf{Ví dụ:} Buổi học bắt đầu lúc $8\text{ giờ } 15\text{ phút}$, kết thúc lúc $9\text{ giờ } 45\text{ phút}$.\\
  Thời gian học: $9\text{h}45 - 8\text{h}15 = \mathbf{1\text{ giờ } 30\text{ phút}}$.
  \item $1\text{ giờ} = 60\text{ phút}$; $1\text{ ngày} = 24\text{ giờ}$; $1\text{ tuần} = 7\text{ ngày}$.
\end{itemize}
}
\end{tcolorbox}
\end{minipage}

\vspace{-2.5mm}
% --- KHỐI 3: SƠ ĐỒ TƯ DUY TIKZ ---
\noindent
\begin{tcolorbox}[
  enhanced,
  colback=orange!4!white,
  colframe=orange!80!black,
  arc=2mm, boxrule=0.8pt,
  left=2mm, right=2mm, top=1.5mm, bottom=1.5mm,
  title={\footnotesize\bfseries\color{white}\faSitemap\ SƠ ĐỒ TƯ DUY: HỆ THỐNG THỜI GIAN VÀ LỊCH LỚP 3}
]
\centering
\begin{tikzpicture}[
  box/.style={draw, rounded corners=2pt, align=center, font=\scriptsize\bfseries, inner sep=2.5pt, minimum height=0.75cm},
  arr/.style={-{Stealth[length=2mm]}, thick}
]
  \node[box, fill=purple!15, draw=purple!80!black, text width=2.6cm] (t1) at (0, 0) {1 GIỜ = 60 PHÚT\\Kim giờ quay 1 số\\Kim phút quay 1 vòng};
  \node[box, fill=teal!15, draw=teal!80!black, text width=2.6cm] (t2) at (3.5, 0) {1 NGÀY = 24 GIỜ\\Từ 0 giờ đến 24 giờ\\Sáng, trưa, chiều, tối};
  \node[box, fill=blue!12, draw=blue!80!black, text width=2.6cm] (t3) at (7.0, 0) {1 TUẦN = 7 NGÀY\\Thứ Hai $\to$ Chủ nhật\\Cách nhau 7 ngày};
  \node[box, fill=orange!15, draw=orange!80!black, text width=2.6cm] (t4) at (10.5, 0) {1 NĂM = 12 THÁNG\\$365$ hoặc $366$ ngày\\Bốn mùa Xuân Hạ Thu Đông};

  \draw[arr, purple!80!black] (t1) -- (t2);
  \draw[arr, teal!80!black] (t2) -- (t3);
  \draw[arr, blue!80!black] (t3) -- (t4);
\end{tikzpicture}
\end{tcolorbox}

\vspace{-2.5mm}
% --- KHỐI 4: VÍ DỤ & MẸO TRÁNH BẪY ---
\noindent
\begin{minipage}[t]{0.49\textwidth}\vspace{0pt}
\begin{tcolorbox}[
  enhanced,
  colback=green!4!white,
  colframe=green!60!black,
  arc=2mm, boxrule=0.8pt,
  left=2mm, right=2mm, top=1.8mm, bottom=1.8mm,
  title={\footnotesize\bfseries\color{white}\faCheckCircle\ Ví dụ thực hành chuẩn mực}
]
{\footnotesize
\textbf{1. Tìm ngày trong tháng:}\\
Nếu ngày $5$ tháng $3$ là ngày \textbf{Thứ Sáu}, thì các ngày Thứ Sáu tiếp theo trong tháng $3$ là:\\
-- Thứ Sáu tuần sau: $5 + 7 = \mathbf{12}$.\\
-- Thứ Sáu tiếp nữa: $12 + 7 = \mathbf{19}$.\\
-- Thứ Sáu cuối cùng: $19 + 7 = \mathbf{26}$.\\
\textbf{2. Chuyển sang tháng tiếp theo:}\\
Ngày $31$ tháng $5$ là Thứ Hai thì ngày $1$ tháng $6$ là \textbf{Thứ Ba} (vì tháng 5 có đúng 31 ngày).
}
\end{tcolorbox}
\end{minipage}
\hfill
\begin{minipage}[t]{0.49\textwidth}\vspace{0pt}
\begin{tcolorbox}[
  enhanced,
  colback=red!4!white,
  colframe=red!70!black,
  arc=2mm, boxrule=0.8pt,
  left=2mm, right=2mm, top=1.8mm, bottom=1.8mm,
  title={\footnotesize\bfseries\color{white}\faExclamationTriangle\ Mẹo tránh bẫy trắc nghiệm}
]
{\footnotesize
\begin{itemize}\setlength{\itemsep}{1.5pt}\setlength{\leftskip}{-2.5mm}
  \item \textbf{Bẫy quên số ngày của tháng khi tính ngày đầu tháng sau:}\\
  Nếu tháng có $30$ ngày (như tháng 4), ngày cuối cùng là 30. Nếu bạn nhớ nhầm thành 31 thì sẽ tính sai toàn bộ thứ của tháng sau!
  \item Luôn nhẩm lại quy tắc nắm tay trước khi tính chuyển tháng.
\end{itemize}
}
\end{tcolorbox}
\end{minipage}


\end{document}
